% Options for packages loaded elsewhere
\PassOptionsToPackage{unicode}{hyperref}
\PassOptionsToPackage{hyphens}{url}
\PassOptionsToPackage{dvipsnames,svgnames,x11names}{xcolor}
%
\documentclass[
  10pt,
  a4paper,
]{article}
\usepackage{amsmath,amssymb}
\usepackage{iftex}
\ifPDFTeX
  \usepackage[T1]{fontenc}
  \usepackage[utf8]{inputenc}
  \usepackage{textcomp} % provide euro and other symbols
\else % if luatex or xetex
  \usepackage{unicode-math} % this also loads fontspec
  \defaultfontfeatures{Scale=MatchLowercase}
  \defaultfontfeatures[\rmfamily]{Ligatures=TeX,Scale=1}
\fi
\usepackage{lmodern}
\ifPDFTeX\else
  % xetex/luatex font selection
\fi
% Use upquote if available, for straight quotes in verbatim environments
\IfFileExists{upquote.sty}{\usepackage{upquote}}{}
\IfFileExists{microtype.sty}{% use microtype if available
  \usepackage[]{microtype}
  \UseMicrotypeSet[protrusion]{basicmath} % disable protrusion for tt fonts
}{}
\makeatletter
\@ifundefined{KOMAClassName}{% if non-KOMA class
  \IfFileExists{parskip.sty}{%
    \usepackage{parskip}
  }{% else
    \setlength{\parindent}{0pt}
    \setlength{\parskip}{6pt plus 2pt minus 1pt}}
}{% if KOMA class
  \KOMAoptions{parskip=half}}
\makeatother
\usepackage{xcolor}
\usepackage[margin=1in]{geometry}
\usepackage{color}
\usepackage{fancyvrb}
\newcommand{\VerbBar}{|}
\newcommand{\VERB}{\Verb[commandchars=\\\{\}]}
\DefineVerbatimEnvironment{Highlighting}{Verbatim}{commandchars=\\\{\}}
% Add ',fontsize=\small' for more characters per line
\newenvironment{Shaded}{}{}
\newcommand{\AlertTok}[1]{\textcolor[rgb]{1.00,0.00,0.00}{\textbf{#1}}}
\newcommand{\AnnotationTok}[1]{\textcolor[rgb]{0.38,0.63,0.69}{\textbf{\textit{#1}}}}
\newcommand{\AttributeTok}[1]{\textcolor[rgb]{0.49,0.56,0.16}{#1}}
\newcommand{\BaseNTok}[1]{\textcolor[rgb]{0.25,0.63,0.44}{#1}}
\newcommand{\BuiltInTok}[1]{\textcolor[rgb]{0.00,0.50,0.00}{#1}}
\newcommand{\CharTok}[1]{\textcolor[rgb]{0.25,0.44,0.63}{#1}}
\newcommand{\CommentTok}[1]{\textcolor[rgb]{0.38,0.63,0.69}{\textit{#1}}}
\newcommand{\CommentVarTok}[1]{\textcolor[rgb]{0.38,0.63,0.69}{\textbf{\textit{#1}}}}
\newcommand{\ConstantTok}[1]{\textcolor[rgb]{0.53,0.00,0.00}{#1}}
\newcommand{\ControlFlowTok}[1]{\textcolor[rgb]{0.00,0.44,0.13}{\textbf{#1}}}
\newcommand{\DataTypeTok}[1]{\textcolor[rgb]{0.56,0.13,0.00}{#1}}
\newcommand{\DecValTok}[1]{\textcolor[rgb]{0.25,0.63,0.44}{#1}}
\newcommand{\DocumentationTok}[1]{\textcolor[rgb]{0.73,0.13,0.13}{\textit{#1}}}
\newcommand{\ErrorTok}[1]{\textcolor[rgb]{1.00,0.00,0.00}{\textbf{#1}}}
\newcommand{\ExtensionTok}[1]{#1}
\newcommand{\FloatTok}[1]{\textcolor[rgb]{0.25,0.63,0.44}{#1}}
\newcommand{\FunctionTok}[1]{\textcolor[rgb]{0.02,0.16,0.49}{#1}}
\newcommand{\ImportTok}[1]{\textcolor[rgb]{0.00,0.50,0.00}{\textbf{#1}}}
\newcommand{\InformationTok}[1]{\textcolor[rgb]{0.38,0.63,0.69}{\textbf{\textit{#1}}}}
\newcommand{\KeywordTok}[1]{\textcolor[rgb]{0.00,0.44,0.13}{\textbf{#1}}}
\newcommand{\NormalTok}[1]{#1}
\newcommand{\OperatorTok}[1]{\textcolor[rgb]{0.40,0.40,0.40}{#1}}
\newcommand{\OtherTok}[1]{\textcolor[rgb]{0.00,0.44,0.13}{#1}}
\newcommand{\PreprocessorTok}[1]{\textcolor[rgb]{0.74,0.48,0.00}{#1}}
\newcommand{\RegionMarkerTok}[1]{#1}
\newcommand{\SpecialCharTok}[1]{\textcolor[rgb]{0.25,0.44,0.63}{#1}}
\newcommand{\SpecialStringTok}[1]{\textcolor[rgb]{0.73,0.40,0.53}{#1}}
\newcommand{\StringTok}[1]{\textcolor[rgb]{0.25,0.44,0.63}{#1}}
\newcommand{\VariableTok}[1]{\textcolor[rgb]{0.10,0.09,0.49}{#1}}
\newcommand{\VerbatimStringTok}[1]{\textcolor[rgb]{0.25,0.44,0.63}{#1}}
\newcommand{\WarningTok}[1]{\textcolor[rgb]{0.38,0.63,0.69}{\textbf{\textit{#1}}}}
\setlength{\emergencystretch}{3em} % prevent overfull lines
\providecommand{\tightlist}{%
  \setlength{\itemsep}{0pt}\setlength{\parskip}{0pt}}
\setcounter{secnumdepth}{-\maxdimen} % remove section numbering
\setlength{\emergencystretch}{3em}
\ifLuaTeX
  \usepackage{selnolig}  % disable illegal ligatures
\fi
\usepackage{bookmark}
\IfFileExists{xurl.sty}{\usepackage{xurl}}{} % add URL line breaks if available
\urlstyle{same}
\hypersetup{
  colorlinks=true,
  linkcolor={blue},
  filecolor={Maroon},
  citecolor={Blue},
  urlcolor={blue},
  pdfcreator={LaTeX via pandoc}}

\author{}
\date{}

\begin{document}

\section{BA1-T3 - Minimal BAE ontology closure
review}\label{ba1-t3---minimal-bae-ontology-closure-review}

\textbf{Revision:} R1

\textbf{Status:} CLOSED / BA1 COMPLETE

\textbf{Repository baseline reviewed:}
\texttt{e88d7e220536863d564f9e3b9fac7f1592a8c440}

\textbf{Phase:} BA1 - Minimal BAE ontology

\textbf{BA0:} CLOSED

\textbf{BA2:} NOT STARTED BY THIS REVIEW

\subsection{1. Closure question}\label{closure-question}

BA1-T3 asks only:

\begin{quote}
Is the refined \texttt{BAReferent\ +\ BAProposition} identity boundary
necessary, sufficient and correctly phased to close BA1 without
importing relation vocabulary, domain taxonomies, document types,
threat-method categories or implementation structure?
\end{quote}

This is a closure and falsification review. It is not another ontology
invention exercise.

The test uses the exit criteria from Work Plan R7: necessity,
sufficiency, no leakage, correct phase placement, the T2 minimality
refinement and a stop criterion against unnecessary BA1-T4 work.

\subsection{2. What ``base entities'' means in
BA1}\label{what-base-entities-means-in-ba1}

The phrase ``base entities that describe the project'' is directionally
useful but needs one distinction.

\texttt{BAReferent} is the family for \textbf{what shared project
meaning is being talked about}. That meaning may be structural,
informational, behavioral, contractual, stateful or otherwise
conceptual; it is not necessarily one runtime object instance.

\texttt{BAProposition} is the family for \textbf{what Base Analysis
asserts about those referents}.

For example, a bounded facial-access projection can distinguish
referents such as:

\begin{itemize}
\tightlist
\item
  CameraSubsystem;
\item
  RecognitionProcessor;
\item
  RecognitionCapture;
\item
  RecognitionRequest;
\item
  the delivery behavior itself when it needs reusable identity;
\item
  the current Ethernet realization if analysis needs to identify it
  independently.
\end{itemize}

Separate propositions can then state source, destination, information
carried, correlation, failure semantics, responsibility/externality and
applicable security constraints.

This avoids forcing every useful noun or behavior into its own metaclass
while preserving enough identity for later relations, provenance and
projections.

\subsection{3. Evidence inherited from BA1-T1 and
BA1-T2}\label{evidence-inherited-from-ba1-t1-and-ba1-t2}

BA1-T1 established the lower-bound distinction between referent identity
and assertion identity. It rejected one fully undifferentiated
analytical element and deliberately left conventional domain roots
unaccepted.

BA1-T2 then attacked the candidate in both directions:

\begin{itemize}
\tightlist
\item
  split pressure across behavior/event, information/resource and
  participant/capability;
\item
  separate split pressure across boundary/externality, store, contract
  and state/context;
\item
  bounded method-consumer pressure using the earlier STRIDE transfer
  probe;
\item
  collapse pressure attempting to remove the semantic distinction
  between BAReferent and BAProposition.
\end{itemize}

No tested semantic kind forced a dedicated first-class split. The
semantic one-family collapse failed. T2 also removed an R1
overextension: project-semantic qualification should target the
underlying BAReferent, not treat a BAProposition record as the project
meaning being constrained.

\subsection{\texorpdfstring{4. Necessity test A - remove
\texttt{BAReferent}}{4. Necessity test A - remove BAReferent}}\label{necessity-test-a---remove-bareferent}

\subsubsection{Hypothesis}\label{hypothesis}

Represent only analytical assertions, using source strings, document
identifiers or assertion-local values for the things those assertions
mention.

\subsubsection{Failure}\label{failure}

The same project meaning is reused across multiple assertions,
projections and baselines. Facial access needs stable recognition of
CameraSubsystem, RecognitionProcessor, RecognitionCapture and
RecognitionRequest across delivery, security and mutation reasoning.
Order fulfillment reuses OrderEvaluation, Reservation, provider/contract
meanings and the physical handoff milestone across branches.

Without referent identity, each proposition or projection must
reconstruct equivalence independently or rely on document syntax as
semantic identity. That breaks or weakens:

\begin{itemize}
\tightlist
\item
  BA0-C2 baseline-scoped shared semantic identity;
\item
  BA0-C5 projection reuse and source drill-down;
\item
  BA0-C6 change-impact traceability;
\item
  the method-neutral reuse required by the T2 consumer probe.
\end{itemize}

\subsubsection{Disposition}\label{disposition}

\begin{Shaded}
\begin{Highlighting}[]
\NormalTok{Remove BAReferent: FAIL}
\NormalTok{BAReferent necessity: SUPPORTED}
\end{Highlighting}
\end{Shaded}

\subsection{\texorpdfstring{5. Necessity test B - remove
\texttt{BAProposition}}{5. Necessity test B - remove BAProposition}}\label{necessity-test-b---remove-baproposition}

\subsubsection{Hypothesis}\label{hypothesis-1}

Keep only BAReferent and attach every fact as anonymous
properties/edges, or reify each assertion as another generic referent.

\subsubsection{Failure}\label{failure-1}

Current evidence requires claim-level origin and evolution distinct from
referent identity. A CameraSubsystem referent may remain stable while a
delivery or realization assertion is retained, revised or retired. One
OrderEvaluation participates in independently evolving claims about
reservation, payment, compensation and commitment.

Anonymous properties/edges can be given IDs and metadata, but once they
have independent assertion identity, origin, diagnosis and change
disposition they implement proposition semantics under another
representation. Reifying every assertion as an undifferentiated referent
likewise hides rather than removes the distinction.

Removing BAProposition therefore undermines:

\begin{itemize}
\tightlist
\item
  BA0-C1 origin/provenance boundary;
\item
  BA0-C4 diagnostic localization;
\item
  BA0-C6 change-impact traceability;
\item
  BA0-C7 source-localized feedback handoff.
\end{itemize}

\subsubsection{Disposition}\label{disposition-1}

\begin{Shaded}
\begin{Highlighting}[]
\NormalTok{Remove BAProposition: FAIL}
\NormalTok{BAProposition necessity: SUPPORTED}
\end{Highlighting}
\end{Shaded}

\subsection{6. Sufficiency test - does a third family remain
hidden?}\label{sufficiency-test---does-a-third-family-remain-hidden}

The closure review rechecks every split pressure raised by R2.

\subsubsection{Behavior / event}\label{behavior-event}

The facial-access delivery and order-fulfillment physical handoff need
reusable semantic identity in some projections. That is sufficient
pressure for BAReferent identity, but not for a universal Behavior/Event
root with closed subtype-specific invariants.

\textbf{Disposition:} NOT FORCED AS FIRST-CLASS SPLIT.

\subsubsection{Information / resource}\label{information-resource}

RecognitionCapture, RecognitionRequest, OrderEvaluation, Reservation and
result concepts need stable identity and lifecycle/correlation
assertions. Current evidence does not establish a common invariant set
that requires a separate Information/Resource identity family.

\textbf{Disposition:} NOT FORCED AS FIRST-CLASS SPLIT.

\subsubsection{Participant / component / capability /
boundary}\label{participant-component-capability-boundary}

Internal and external responsibility placement must be explicit, but it
can be expressed through referent identities plus method-neutral
roles/propositions. Service consumption does not transfer ownership.

\textbf{Disposition:} NOT FORCED AS FIRST-CLASS SPLIT.

\subsubsection{Store / contract}\label{store-contract}

InventoryLedger, ReservationStore, PaymentProviderContract and
CarrierContract can have stable BAReferent identity and distinct
classifications. The reviewed evidence does not require separate root
invariants for all stores or contracts.

\textbf{Disposition:} NOT FORCED AS FIRST-CLASS SPLIT.

\subsubsection{State / context}\label{state-context}

Local states may remain values; asserted conditions/transitions may be
propositions; independently reusable named states or contexts may
receive BAReferent identity. No universal State/Context root is
required.

\textbf{Disposition:} NOT FORCED AS FIRST-CLASS SPLIT.

\subsubsection{Diagnostic / provenance /
view}\label{diagnostic-provenance-view}

These are real BA0 responsibilities but their independent representation
belongs to BA3/BA4 metadata and projection contracts. Promoting them to
domain BAE roots would confuse ontology with analytical administration.

\textbf{Disposition:} CORRECTLY DEFERRED.

\subsubsection{Sufficiency result}\label{sufficiency-result}

No material identity responsibility remains unrepresented. The two
families are sufficient at the BA1 abstraction level.

\subsection{7. No-leakage review}\label{no-leakage-review}

The closure boundary does not require:

\begin{itemize}
\tightlist
\item
  MR/Decision/FR/SR document kinds as BAE types;
\item
  STRIDE, STRIDE-AI or DFD roots;
\item
  actor/component/asset/boundary/data\_flow historical taxonomy;
\item
  SysML/KerML/ArchiMate/AADL metaclasses;
\item
  ThreatForge classes;
\item
  one graph, database, programming class hierarchy or diagram notation;
\item
  a fixed semantic-kind vocabulary.
\end{itemize}

This preserves BA0-C3 method neutrality and BA0-C8 representation
independence.

\subsection{8. Specialization review - types versus
classifications}\label{specialization-review---types-versus-classifications}

BA1 closure explicitly distinguishes three levels:

\begin{Shaded}
\begin{Highlighting}[]
\NormalTok{first{-}class identity family}
\NormalTok{    BAReferent / BAProposition}

\NormalTok{method{-}neutral semantic kind or role}
\NormalTok{    behavior, information, capability, contract, store, state, ...}

\NormalTok{method{-}specific interpretation}
\NormalTok{    STRIDE/DFD/etc. categories owned by the analysis method}
\end{Highlighting}
\end{Shaded}

The middle level is important: the project can distinguish a delivery
behavior from information or a contract without requiring three
independent ontology roots.

BA2 may define how proposition participation roles, relation/action
semantics and reusable semantic classifications are represented. BA5 may
later address controlled lexical vocabulary and authoring assistance.
Neither phase may silently promote a category into a new BA1 root
without a concrete reopen trigger.

\subsection{9. Correct phase-placement
review}\label{correct-phase-placement-review}

The unresolved questions after BA1 are real, but they do not block
identity-ontology closure.

\begin{itemize}
\tightlist
\item
  proposition arity, participant roles, relation/action semantics
  -\textgreater{} \textbf{BA2};
\item
  identity/equivalence mechanics, source locators,
  grounded/derived/diagnostic materialization, stale/change state
  -\textgreater{} \textbf{BA3};
\item
  human/method views and projection contracts -\textgreater{}
  \textbf{BA4};
\item
  controlled vocabulary, lexical normalization and optional assistance
  -\textgreater{} \textbf{BA5};
\item
  full Base Analysis regression and implementation-independent closure
  -\textgreater{} \textbf{BA6};
\item
  AnalysisRecord, Finding and formal method overlays -\textgreater{}
  \textbf{after BA6}.
\end{itemize}

No unresolved BA1 identity question is being hidden inside BA2-BA6.

\subsection{10. Minimality refinement
review}\label{minimality-refinement-review}

T2 removed proposition-as-project-target from the required core. BA1-T3
confirms that reduction.

If a delivery, milestone, policy-like meaning or named state must be
qualified by multiple project-semantic assertions, that meaning receives
BAReferent identity. BAProposition remains the assertion record carrying
claim-level provenance/diagnostic/change pressure.

This avoids confusing ``the thing being constrained'' with ``the
analytical statement describing the constraint''.

\subsection{11. First-class split
criterion}\label{first-class-split-criterion}

BA1-T3 closes the following criterion:

\begin{quote}
A new dedicated BAE type/family requires both independent semantic
identity across propositions/projections/change and reusable
subtype-specific invariants that cannot be represented honestly as
classifications, roles, values or propositions over the accepted
families.
\end{quote}

This criterion is deliberately stronger than ``the concept appears
often'' or ``a tool/method has a class for it''.

\subsection{12. Reopen review}\label{reopen-review}

Future specialization is allowed, but not by silent ontology growth.

BA1 must be reopened if later concrete evidence proves, for example,
that a behavior/event, information, state, boundary, contract or other
semantic kind needs subtype-specific invariants that generic BAReferent
plus BAProposition cannot preserve for method-neutral consumers.

Ordinary implementation convenience, prettier diagrams or one analysis
method's preferred taxonomy do not reopen BA1.

\subsection{13. Stop criterion}\label{stop-criterion}

The closure review found:

\begin{itemize}
\tightlist
\item
  both accepted families are necessary;
\item
  no tested third identity family is necessary;
\item
  the two-family boundary is method-neutral and
  representation-independent;
\item
  residual questions have valid owners in BA2-BA6;
\item
  the T2 reduction removes rather than adds ontology surface.
\end{itemize}

A BA1-T4 experiment would therefore be reassurance work without a
specific unresolved falsification target.

\subsubsection{Disposition}\label{disposition-2}

\begin{Shaded}
\begin{Highlighting}[]
\NormalTok{BA1{-}T4: NOT REQUIRED}
\end{Highlighting}
\end{Shaded}

\subsection{14. BA1 closure decision}\label{ba1-closure-decision}

BA1 is closed for the current thesis scope.

\begin{Shaded}
\begin{Highlighting}[]
\NormalTok{BA1                             CLOSED}
\NormalTok{BAReferent                      ACCEPTED}
\NormalTok{BAProposition                   ACCEPTED}
\NormalTok{BAE umbrella metaclass          NOT REQUIRED}
\NormalTok{Dedicated domain root splits    NOT REQUIRED BY CURRENT EVIDENCE}
\NormalTok{Proposition{-}as{-}project{-}target   NOT REQUIRED}
\NormalTok{BA2                             NOT STARTED / NEXT AUTHORIZED PHASE}
\end{Highlighting}
\end{Shaded}

The accepted boundary is recorded canonically in
\texttt{BA1\_MINIMAL\_BAE\_IDENTITY\_ONTOLOGY\_R1.md}.

\subsection{15. Next authorized
microstep}\label{next-authorized-microstep}

After this closure package is reviewed, committed and remotely verified,
execute only:

\begin{quote}
\textbf{BA2-T1 - minimal BAProposition structural shape and
participation-role lower-bound derivation.}
\end{quote}

BA2-T1 must ask what minimum methodology-neutral structure allows a
BAProposition to state reusable facts over BAReferents without yet
freezing an exhaustive relation/action vocabulary or importing
method-specific semantics.

\end{document}
